\section{Related Work}
\label{sec:related_work}

Using Large Language Models (LLMs) to automate unit test generation offers a compelling alternative to traditional random or search-based testing. We structure our review of recent progress in this field around two main areas: zero-shot LLM test generation and feedback-driven iterative refinement. We then highlight the gaps in these existing works and position our study accordingly.

\subsection{LLM-based Zero-shot Unit Test Generation}
\label{subsec:zero_shot_gen}

Initial work in LLM-based testing focused on whether models could generate working unit tests immediately, without any iterative fixing. Guilherme and Vincenzi~\cite{guilherme2023sast} took this approach with \texttt{gpt-3.5-turbo}, evaluating it on 33 Java projects. Under a 0.2 temperature and carefully formatted prompts, the model achieved a 64.6\% build success rate, 90.2\% average coverage, and a 70.5\% mutation score. Sch{\"a}fer et al.~\cite{schafer2024testpilot} followed up with TestPilot, a tool they used to benchmark \texttt{gpt-3.5-turbo}, \texttt{code-cushman-002}, and StarCoder across JavaScript, TypeScript, and Python. While they reported a median statement coverage of 70.2\% and branch coverage of 52.8\%, the resulting test suites often failed to compile due to syntax issues or broken external dependencies.

Looking beyond single languages, Celik and Mahmoud~\cite{celik2026gem} developed the GEM framework to benchmark Codex, Sonnet, and Qwen on Java, Python, and C++ code. On the ULT(100) dataset, Codex hit 84.30\% branch coverage and a 72.71\% mutation score, though the authors noticed significant performance gaps when switching languages. In a different approach, Baskaran et al.~\cite{baskaran2025jucct} turned to multi-agent architectures, orchestrating OpenAI models with AutoGen to test a Python insurance app. This team-of-agents setup pushed branch coverage to 99\% and mutation scores to 95.8\%, though they only evaluated the framework on one specific software system.

However, zero-shot test generation struggles heavily with compilation and execution rates on the first try. Lops et al.~\cite{lops2025agonetest} evaluated newer models like \texttt{gpt-4o-mini} and \texttt{gemini-1.5-pro} on the class-level Java benchmark CLASSES2TEST. Although the models reached up to 79.8\% branch coverage and 89.2\% mutation scores, the compilation success rate peaked at just 36.0\%, showing a massive divide between test quality metrics and actual runnability. Crucially, these existing benchmarks do not control for the structural characteristics of the input code, such as its cyclomatic complexity. As a result, we still lack a clear understanding of how code complexity affects the build rates and quality of initial LLM outputs.

\subsection{Feedback-Driven and Iterative Test Refinement}
\label{subsec:iterative_refinement}

Recognizing that first-try tests often fail, researchers began adding compiler and execution feedback to let models debug their own code. ChatTester, proposed by Yuan et al.~\cite{yuan2024chattester}, is one such tool; it feeds compilation errors back to the LLM to repair Java class tests. By letting the model learn from its mistakes, ChatTester boosted compilation rates by 34.3\%, eventually reaching an 82\% branch coverage and a 65\% mutation score.

Along a similar line, Pan et al.~\cite{pan2024aster} designed ASTER, which pairs GPT-4 with static analysis and context extraction to resolve errors in Java and Python tests. ASTER reached 78.0\% line coverage and 77.2\% branch coverage on Python benchmarks, winning approval in a developer study. Instead of fixing compiler errors, Straubinger et al.~\cite{straubinger2025mutation} used mutation testing results to drive a debugging loop, raising branch coverage from 55\% to nearly 80\% and mutation scores from 60\% to 80\%. Lu et al.~\cite{lu2026beyond} took this further by using DeepSeek-V3 to target survived mutants specifically, improving mutation scores by 16.04\% for standalone methods and 8.11\% for non-standalone methods.

Even though these feedback loops improve coverage and compilability, they come with high operational costs. Multiple round-trips to LLM APIs inflate token costs, increase execution delays, and require complex local runtimes and static analysis setups. More importantly, because these frameworks combine generation and healing, they hide the base model's standalone zero-shot capabilities. We cannot easily see how source code complexity impacts the initial output quality when it is immediately masked by corrective iterations.

\subsection{Research Positioning}
\label{subsec:positioning}

We position our work to address two specific gaps. First, past evaluations of zero-shot generation do not isolate cyclomatic complexity; they use benchmarks with arbitrary complexity, making it hard to see if complex code is what breaks the LLM's performance. We address this by testing \texttt{gpt-4o-mini-2024-07-18} on a balanced, complexity-controlled dataset of Java and Python methods (with cyclomatic complexity strictly between 5 and 15). Second, most current papers rely on simple descriptive statistics to compare LLMs with random testing baselines, lacking statistical rigor. We resolve this by applying non-parametric tests---specifically the one-sided sign test per language and the paired Wilcoxon signed-rank test---to confirm whether the improvements in branch coverage and mutation scores are statistically significant under a pre-registered analysis plan.
